
\section{Composite Problem via Duality}


\subsection{Motivation}

In this section, we will discuss a few algorithmic applications of
duality. In general, two reasons to consider a dual problem are that
(1) the dual problem might have smaller dimension and (2) (sub)gradients
are easier to compute. We will see such examples. 

Suppose we have a ``difficult'' convex problem $\min_{x}f(x)$,
i.e., one that does not satisfy some smoothness condition which direct
suggests a faster than worst-case optimization method. Often, we can
decompose the difficult problem as $f(x)=g(x)+h(Ax)$ such that the
(sub)gradients $\nabla g^{*}(x)$ and $\nabla h^{*}(x)$ are easy
to compute. Then, it is useful to compute its dual as follows:
\begin{align*}
\min_{x}g(x)+h(Ax) & =\min_{x}\max_{\theta}g(x)+\theta^{\top}Ax-h^{*}(\theta)\\
 & =\max_{\theta}\min_{x}g(x)+(A^{\top}\theta)^{\top}x-h^{*}(\theta)\\
 & =\max_{\theta}-g^{*}(-A^{\top}\theta)-h^{*}(\theta)\\
 & =-\min_{\theta}g^{*}(-A^{\top}\theta)+h^{*}(\theta)
\end{align*}
where we used $h=h^{**}$ in the first line, the following minimax
theorem on the second line, and the definition of $g^{*}$ on the
third line.
\begin{thm}[Sion's minimax theorem]
Let $X\subset\Rn$ be a compact convex set and $Y\subset\Rm$ be
a convex set. If $f:X\times Y\rightarrow\R\cup\{+\infty\}$ such that
$f(x,\cdot)$ is upper semi-continuous and quasi-concave on $Y$ for
all $x\in X$ and $f(\cdot,y)$ is lower semi-continuous and quasi-convex
on $X$ for all $y\in Y$. Then, we have
\[
\min_{x\in X}\sup_{y\in Y}f(x,y)=\sup_{y\in Y}\min_{x\in X}f(x,y).
\]
\end{thm}

\begin{rem*}
Compactness is necessary. Consider $f(x,y)=x+y$. This theorem generalizes
Von Neumann's minimax theorem. 
\end{rem*}
The above idea of passing to the dual extends to a sum of convex functions
as below below.
\begin{lem}
\label{lem:dual-decomp}For convex functions $f_{1},\dots,f_{m}$
with closed epigraphs, we have (assuming the minimum exists)
\[
\min_{x}\sum_{i=1}^{m}f_{i}(x)=-\min_{\theta_{1},\ldots,\theta_{m-1}}\left(f_{1}^{*}(-\theta_{1})+\sum_{i=2}^{m-1}f_{i}^{*}(\theta_{i-1}-\theta_{i})+f_{m}^{*}(\theta_{m-1})\right)
\]
\end{lem}

\begin{proof}
For any $\theta_{1},\ldots,\theta_{m-1}$ and $x$, 
\begin{align*}
f_{1}^{*}(-\theta_{1})+\sum_{i=2}^{m-1}f_{i}^{*}(\theta_{i-1}-\theta_{i})+f_{m}^{*}(\theta_{m-1}) & \ge((-\theta_{1})^{\top}x-f_{1}(x))+\sum_{i=1}^{m-1}((\theta_{i-1}-\theta_{i})^{\top}x-f_{i}(x))\\
 & ~+(\theta_{m-1}^{\top}x-f_{m}(x))\\
 & =-\sum_{i=1}^{m}f_{i}(x)\\
\sum_{i=1}^{m}f_{i}(x) & \ge-(f_{1}^{*}(-\theta_{1})+\sum_{i=2}^{m-1}f_{i}^{*}(\theta_{i-1}-\theta_{i})+f_{m}^{*}(\theta_{m-1}))
\end{align*}
So, 
\begin{align*}
\inf_{x}\sum_{i=1}^{m}f_{i}(x) & \ge\sup_{\theta_{1},\ldots,\theta_{m-1}}-(f_{1}^{*}(-\theta_{1})+\sum_{i=2}^{m-1}f_{i}^{*}(\theta_{i-1}-\theta_{i})+f_{m}^{*}(\theta_{m-1}))
\end{align*}
It remains to show that there exist $x$ and $\theta_{1},\ldots,\theta_{m-1}$
such that 
\[
\sum_{i=1}^{m}f_{i}(x)\le-(f_{1}^{*}(-\theta_{1})+\sum_{i=2}^{m-1}f_{i}^{*}(\theta_{i-1}-\theta_{i})+f_{m}^{*}(\theta_{m-1}))
\]
($\ge$ is guaranteed by above). We choose $x\in\R^{n}$ to minimize
the convex function $\sum_{i=1}^{m}f_{i}(x)$. Then, $0\in\nabla(\sum_{i=1}^{m}f_{i}(x))=\sum_{i=1}^{m}\nabla f_{i}(x)$
(as Minkowski sum of subgradients), so there exists $\alpha_{1},\ldots,\alpha_{m}\in\R^{n}$
such that $\sum_{i=1}^{m}\alpha_{i}=0$ and $\alpha_{i}\in\nabla f_{i}(x)$.
By definition of subgradient, for all $y\in\R^{n}$, 
\begin{align*}
f_{i}(y) & \ge f_{i}(x)+\alpha_{i}^{\top}(y-x)\\
\alpha_{i}^{\top}y-f_{i}(y) & \le\alpha_{i}^{\top}x-f_{i}(x)\\
f_{i}^{*}(\alpha_{i}) & =\sup_{y}\alpha_{i}^{\top}y-f_{i}(y)\\
 & \le\alpha_{i}^{\top}x-f_{i}(x)
\end{align*}
Then, 
\begin{align*}
\sum_{i=1}^{m}f_{i}^{*}(\alpha_{i}) & \le\sum_{i=1}^{m}\left(\alpha_{i}^{\top}x-f_{i}(x)\right)\\
 & =-\sum_{i=1}^{m}f_{i}(x)
\end{align*}
Finally, we choose $\theta_{i}=-\sum_{j=1}^{i}\alpha_{j}$ for $i=0,\ldots,m$,
so that $\alpha_{i}=\theta_{i-1}-\theta_{i}$, and $\theta_{0}=\theta_{m}=0$.
Then, 
\[
\sum_{i=1}^{m}f_{i}(x)\le-\sum_{i=1}^{m}f_{i}^{*}(\theta_{i-1}-\theta_{i})
\]
as desired.
\end{proof}
\begin{rem*}
The alternative form in Lemma \ref{lem:dual-decomp} can be used to
implement a distributed version of gradient descent. If server $i$
only has access to $f_{i}$, then it suffices to communicate with
only neighboring servers $i\pm1$ to compute the contribution of $\theta_{i}$
to the gradient, as the only terms with $\theta_{i}$ are $f_{i}^{*}(\theta_{i-1}-\theta_{i}),f_{i+1}^{*}(\theta_{i}-\theta_{i+1})$
\end{rem*}
We call ``$g(x)+h(Ax)$'' the primal problem and ``$g^{*}(-A^{\top}\theta)+h^{*}(\theta)$''
the dual problem. Often, the dual problem gives us some insight on
the primal problem. However, we note that there are many ways to split
a problem into two and hence many candidate dual problems.
\begin{example}
Consider the unit capacity flow problem on a graph $G=(V,E)$:
\[
\max_{Af=d,-1\leq f\leq1}c^{\top}f
\]
where $f\in\R^{E}$ is the flow vector, $A\in\R^{V\times E}$ encodes
the vertex-edge adjacency matrix with two nonzeros per column, $d$
is the demand vector so that $Af=d$ is flow conservation, and $c$
is the cost vector. We can write the dual as follows:
\begin{align*}
\max_{Af=d,-1\leq f\leq1}c^{\top}f & =\max_{-1\leq f\leq1}\min_{\phi}c^{\top}f-\phi^{\top}(Af-d)\\
 & =\min_{\phi}\max_{-1\leq f\leq1}\phi^{\top}d+(c-A^{\top}\phi)^{\top}f\\
 & =\min_{\phi}\phi^{\top}d+\sum_{e\in E}|c-A^{\top}\phi|_{e}.
\end{align*}
\end{example}

When $c=0$ and $d=F\cdot1_{st}$ this is the maximum flow problem
with flow value $F$, and the dual problem is the minimum $s-t$ cut
problem with the cut given by $\{v\in V\text{ such that }\phi(v)\geq t\}$.
We can view $\phi$ as assigning a ``potential'' to every vertex
of the graph. Note that there are $|E|$ variables in primal and $|V|$
variables in dual. So, in this sense, the dual problem is easier for
dense graphs. Although we do not have a way to turn a minimum $s-t$
cut to a maximum $s-t$ flow in general, we will see various tricks
to reconstruct the primal solution from the dual solution by modifying
the problem.

\subsection{Example: Semidefinite Programming}

When you can solve the dual problem, the proof of the optimality of
the dual problem often directly gives you the solution of the primal
problem, or gives you some way to solve the primal problem efficiently.
Here, we use the semidefinite programming as an concrete example.
Define $\bullet$ as the inner product of matrices, i.e., 
\[
A\bullet B=\sum_{i,j}A_{ij}B_{ij}.
\]
Consider the semidefinite programming (SDP) problem:
\begin{equation}
\max_{X\succeq0}C\bullet X\text{ s.t. }A_{i}\bullet X=b_{i}\text{ for }i=1,2,\cdots,m\label{eq:SDP}
\end{equation}
and its dual
\begin{equation}
\min_{y}b^{\top}y\text{ s.t. }\sum_{i=1}^{m}y_{i}A_{i}\succeq C\label{eq:dual_SDP}
\end{equation}
where $X$, $C$, $A_{i}$ are $n\times n$ symmetric matrices and
$b,y\in\Rm$. SDP is a generalization of linear programming and is
useful for various problems involving matrices. If we apply the current-best
cutting plane method naively on the primal problem, we would get $\tilde{O}(n^{2}(Z+n^{4}))$
time algorithm for the primal (because there are $n^{2}$ variables)
and $\tilde{O}(m(Z+n^{\omega}+m^{2}))$ for the dual where $Z$ is
the total number of non-zeros in $A_{i}$. (The term $Z+n^{\omega}$
is the cost of computing $\sum y_{i}A_{i}$ and finding its minimum
eigenvalue.) Generally, $n^{2}\gg m$ and hence it takes much less
time to solve the dual.

We note that
\[
\min_{\sum_{i=1}^{m}y_{i}A_{i}\succeq C}b^{\top}y=\min_{v^{\top}(\sum_{i=1}^{m}y_{i}A_{i}-C)v\geq0\ \forall\norm v_{2}=1}b^{\top}y.
\]
In each step of the cutting plane method, the (sub)gradient oracle
either outputs $b$ or outputs one of the cutting planes
\[
v^{\top}(\sum_{i=1}^{m}y_{i}A_{i}-C)v\geq0.
\]
Let $S$ be the set of all cutting planes used in the algorithm. Then,
the proof of the cutting plane method shows that
\begin{equation}
\min_{\sum_{i=1}^{m}y_{i}A_{i}\succeq C}b^{\top}y=\min_{v^{\top}(\sum_{i=1}^{m}y_{i}A_{i}-C)v\geq0\ \forall v\in\mathcal{S}}b^{\top}y\pm\varepsilon.\label{eq:SDP_apx}
\end{equation}
The key idea to obtaining the primal solution is to take the dual
of the right hand side (which is an approximate dual of the original
problem). Now, we have
\begin{align*}
\min_{v^{\top}(\sum_{i=1}^{m}y_{i}A_{i}-C)v\geq0\ \forall v\in\mathcal{S}}b^{\top}y & =\min_{y}\max_{\lambda_{v}\geq0}b^{\top}y-\sum_{v\in\mathcal{S}}\lambda_{v}v^{\top}(\sum_{i=1}^{m}y_{i}A_{i}-C)v\\
 & =\max_{\lambda_{v}\geq0}\min_{y}C\bullet\sum_{v\in\mathcal{S}}\lambda_{v}vv^{\top}+b^{\top}y-\sum_{i=1}^{m}y_{i}\sum_{v\in\mathcal{S}}\lambda_{v}vv^{\top}\bullet A_{i}\\
 & =\max_{X=\sum_{v\in\mathcal{S}}\lambda_{v}vv^{\top},\lambda_{v}\geq0}\min_{y}C\bullet X+\sum_{i=1}^{m}y_{i}(b_{i}-X\bullet A_{i})\\
 & =\max_{X=\sum_{v\in\mathcal{S}}\lambda_{v}vv^{\top},\lambda_{v}\geq0,X\bullet A_{i}=b_{i}}C\bullet X.
\end{align*}
Note that this is exactly the primal SDP problem, except that we restrict
the set of solutions to the form $\sum_{v\in\mathcal{S}}\lambda_{v}vv^{\top}$
with $\lambda_{v}\geq0$. Also, we can write this problem as a linear
program:
\begin{equation}
\max_{\sum_{v}\lambda_{v}v^{\top}A_{i}v=b_{i}\text{ for all }i,\lambda_{v}\geq0}\sum_{v}\lambda_{v}v^{\top}Cv.\label{eq:SDP_LP}
\end{equation}
Therefore, we can simply solve this linear program and recover an
approximate solution for the SDP. By (\ref{eq:SDP_apx}), we know
that this is an approximate solution with the same guarantee as the
dual SDP. 

Now, we analyze the runtime of this algorithm. This algorithm contains
two phases: solve the dual SDP via cutting plane method, and solve
the primal linear program. Note that each step of the cutting plane
method involves finding a separating hyperplane of $\sum_{i=1}^{m}y_{i}A_{i}\succeq C$. 
\begin{xca}
Let $\Omega=\{y\in\Rm:\ \sum_{i=1}^{m}y_{i}A_{i}\succeq C\}$. Show
that one can implement the separation oracle in time $O^{*}(Z+n^{\omega})$
via eigenvalue computation.
\end{xca}

Therefore, the first phase takes $O^{*}(m(Z+n^{\omega}+m^{2}))$ time
in total. Since the cutting plane method takes $O^{*}(m)$ steps,
we have $|\mathcal{S}|=O^{*}(m)$. In the second phase, we need to
solve a linear program (\ref{eq:SDP_LP}) with $O^{*}(m)$ variables
with $O(m)$ constraints. It is known how to solve such linear programs
in time $O^{*}(m^{2.38})$ \cite{cohen2018solving}. Hence, the total
cost is dominated by the first phase 
\[
O^{*}\left(mZ+mn^{\omega}+m^{3}\right).
\]

\begin{problem}
In the first phase, each step involves computing an eigenvector of
similar matrices. So, can we use matrix update formulas to decrease
the cost per step in the cutting plane to $O^{*}(Z+n^{2})$? Namely,
can we solve SDP in time $O^{*}\left(mZ+m^{3}\right)?$
\end{problem}


\subsection{Duality and Convex Hull}

The problem $\min_{x}g(x)+h(Ax)$ in general can be solved by a similar
trick. To make this geometric picture clearer, we consider its linear
optimization version: $\min_{(x,t_{1})\in\epi\,g,(Ax,t_{2})\in\epi\,h}t_{1}+t_{2}.$
To simplify the notation, we consider the problem
\[
\min_{x\in K_{1},Mx\in K_{2}}c^{\top}x
\]
where $M\in\R^{m\times n}$, and we have convex sets $K_{1}\subset\Rn$
and $K_{2}\subset\Rm$. 

To be concrete, let us consider the following example. Let $V_{1}$
be a set of students, $V_{2}$ be a set of schools. Each edge $e\in E$
represents a possible choice of a student. Let $w_{e}$ be the happiness
of a school/student if the student is assigned to that school. Suppose
that every student can only be assigned to one school and school $b$
can accept $c_{b}$ students. Then, the problem can be formulated
as
\[
\max_{x_{e}\geq0}\sum_{e\in E}w_{e}x_{e}\text{ subject to }\sum_{(a,b)\in E}x_{(a,b)}\leq1\ \forall a\in V_{1},\sum_{(a,b)\in E}x_{(a,b)}\leq c_{b}\ \forall b\in V_{2}.
\]
This is the weighted b-matching problem. Typically, the number of
students is much more than the number of schools. Therefore, an algorithm
with running time linear in the number of students is preferable.
To apply our framework, we let
\begin{align*}
K_{1} & =\{x\in\R^{E},x_{e}\geq0,\sum_{(a,b)\in E}x_{(a,b)}\leq1\ \forall a\in V_{1}\},\\
K_{2} & =\{y\in\R^{V_{2}},y_{b}\leq c_{b}\},
\end{align*}
and $M\in\R^{E}\rightarrow\R^{V_{2}}$ is the map $(Tx)_{b}=\sum_{a:(a,b)\in E}x_{(a,b)}.$

To further emphasize its importance, consider some general examples
here:
\begin{itemize}
\item Linear programming: $\min_{Ax=b,x\geq0}c^{\top}x$: $K_{1}=\{x\geq0\}$,
$K_{2}=\{b\}$ and $M=A$.
\item Semidefinite programming: $\min_{A_{i}\bullet X=b_{i},X\succeq0}C\bullet X$:
$K_{1}=\{X\succeq0\}$, $K_{2}=\{b\}$ and $M:\R^{n\times n}\rightarrow\Rm$
defined by $(MX)_{i}=A_{i}\bullet X$.
\item Matroid intersection: $\min_{x\in M_{1}\cap M_{2}}1^{\top}x$: $K_{1}=M_{1}$
and $K_{2}=M_{2}$, $M=I$.
\item Submodular minimization: $K_{1}=\{y\in\Rn:\sum_{i\in S}y_{i}\leq f(S)\text{ for all }S\subset[n]\}$,
$K_{2}=\{y\leq0\}$, $M=I$.
\item Submodular flow: $K_{1}=\{\varphi\in\R^{E},\ell_{e}\leq\varphi_{e}\leq u_{e}\text{ for all }e\in E\}$,
$K_{2}=\{y\in\R^{V}:\sum_{i\in S}y_{i}\leq f(S)\text{ for all }S\subset[n]\}$,
$M$ is the incidence matrix.
\end{itemize}
In all of these examples, it is easy to compute the gradient of $\delta_{K_{1}}^{*}$
and $\delta_{K_{2}}^{*}$. For the last three examples, it is not
clear how to compute gradient of $\delta_{K_{1}}$ and/or $\delta_{K_{2}}$
directly. Furthermore, in all examples, $M$ maps from a larger space
to the same or smaller space. Therefore, it is good to take advantage
of the smaller space.

Before \cite{lee2015faster}, the standard way was to use the equivalence
of $\nabla\delta_{K_{1}}^{*}$ and $\nabla\delta_{K_{1}}$ , and apply
cutting plane methods. With the running time of cutting plane methods,
such an algorithm usually had theoretical running time at least $n^{5}$
and was of little practical value. 

Now we rewrite the problem as we did in the beginning:
\begin{align*}
\min_{x\in K_{1},Mx\in K_{2}}c^{\top}x & =\min_{x\in\Rn}c^{\top}x+\delta_{K_{1}}(x)+\delta_{K_{2}}(Mx)\\
 & =\min_{x\in\Rn}\max_{\theta\in\Rm}c^{\top}x+\delta_{K_{1}}(x)+\theta^{\top}Mx-\delta_{K_{2}}^{*}(\theta)\\
 & =\max_{\theta\in\Rm}\min_{x\in\Rn}c^{\top}x+\delta_{K_{1}}(x)+\theta^{\top}Mx-\delta_{K_{2}}^{*}(\theta)\\
 & =\max_{\theta\in\Rm}-\delta_{K_{1}}^{*}(-c-M^{\top}\theta)-\delta_{K_{2}}^{*}(\theta).
\end{align*}
Taking the dual has two benefits. First, the number of variables is
smaller. Second, the gradient oracle is something we can compute efficiently.
Hence, cutting plane methods can be used to solve it in $O^{*}(m\mathcal{T}+m^{3})$
where $\mathcal{T}$ is the time to evaluate $\nabla\delta_{K_{1}}^{*}$
and $\nabla\delta_{K_{2}}^{*}$. The only problem left is to recover
the primal $x$.
\begin{proof}
\begin{figure}
\centering{}\includegraphics[width=0.5\columnwidth]{Figures/Lemma_4\lyxdot 47}\caption{The oracle behaves identically on $K_{1}$ and $\widetilde{K_{1}}$ }
\end{figure}
\end{proof}
The key observation is the following lemma:
\begin{lem}
Let $x_{i}\in K_{1}$ be the set of points output by the oracle $\nabla\delta_{K_{1}}^{*}$
during the cutting plane method. Define $y_{i}\in K_{2}$ similarly.
Suppose that the cutting plane method ends with the guarantee that
the additive error is less than $\varepsilon$. Then, we have that
\[
\min_{x\in K_{1},Tx\in K_{2}}c^{\top}x\leq\min_{x\in\widetilde{K_{1}},Tx\in\widetilde{K_{2}}}c^{\top}x\leq\min_{x\in K_{1},Tx\in K_{2}}c^{\top}x+\varepsilon
\]
where $\widetilde{K_{1}}=\conv(x_{i})$ and $\widetilde{K_{2}}=\conv(y_{i})$.
\end{lem}

\begin{proof}
Let $\theta_{i}$ be the set of directions queried by the oracle for
$\nabla\delta_{K_{1}}^{*}$ and $\varphi_{i}$ be the directions queried
by the oracle for $\nabla\delta_{K_{2}}^{*}$. We claim that $x_{i}\in\nabla\delta_{\widetilde{K_{1}}}^{*}(\theta_{i})$
and $y_{i}\in\nabla\delta_{\widetilde{K_{2}}}^{*}(\varphi_{i})$.
Having this, the algorithm cannot distinguish between $\widetilde{K_{1}}$
and $K_{1}$, and between $\widetilde{K_{2}}$ and $K_{2}$. Hence,
the algorithm runs exactly the same, i.e., uses the same sequence
of points. Therefore, we get the same value $c^{\top}x$. However,
by the guarantee of cutting plane method, we have that
\[
\min_{x\in\widetilde{K_{1}},Tx\in\widetilde{K_{2}}}c^{\top}x\leq c^{\top}x\leq\min_{x\in K_{1},Tx\in K_{2}}c^{\top}x+\varepsilon.
\]

To prove the claim, we note that $x_{i}\in\nabla\delta_{\widetilde{K_{1}}}^{*}(\theta_{i})$.
Note that $\widetilde{K_{1}}\subset K_{1}$ and hence $\min_{x\in\widetilde{K_{1}}}\theta_{i}^{\top}x\geq\min_{x\in K_{1}}\theta_{i}^{\top}x$.
Also, note that 
\[
x_{i}=\arg\min_{x\in K_{1}}\theta_{i}^{\top}x\in\widetilde{K_{1}}
\]
and hence $\min_{x\in\widetilde{K_{1}}}\theta_{i}^{\top}x\leq\min_{x\in K_{1}}\theta_{i}^{\top}x$.
Therefore, we have that $\min_{x\in\widetilde{K_{1}}}\theta_{i}^{\top}x=\min_{x\in K_{1}}\theta_{i}^{\top}x$.
Therefore, $x_{i}\in\arg\min_{x\in K_{1}}\theta_{i}^{\top}x$. This
proves the claim for $\widetilde{K_{1}}$. The proof for $\widetilde{K_{2}}$
is the same.
\end{proof}
This reduces the problem into the form $\min_{x\in\widetilde{K_{1}},Tx\in\widetilde{K_{2}}}c^{\top}x$.
For the second phase, we let $z_{i}=Mx_{i}\in\Rm$. Then, we have
\begin{align*}
\min_{x\in\widetilde{K_{1}},Mx\in\widetilde{K_{2}}}c^{\top}x= & \min_{t_{i}\geq0,s_{i}\geq0,M\sum_{i}t_{i}x_{i}=\sum_{i}s_{i}y_{i}}c^{\top}(\sum_{i}t_{i}x_{i})\\
= & \min_{t_{i}\geq0,s_{i}\geq0,\sum_{i}t_{i}z_{i}=\sum_{i}s_{i}y_{i}}\sum t_{i}\cdot c^{\top}x_{i}.
\end{align*}
Note that it takes $O^{*}(mZ)$ time to write down this linear program
where $Z$ is the number of non-zeros in $M$. Next, we note that
this linear program has $O^{*}(m)$ variables and $m$ constraints.
Therefore, we can solve it in $O^{*}(m^{2.38})$ time. 

Therefore, the total running time is 
\[
O^{*}(m(\mathcal{T}+m^{2})+(mZ+m^{2.38})).
\]

To conclude, we have the following theorem.
\begin{thm}
Given convex sets $K_{1}\subset\Rn$ and $K_{2}\subset\Rm$ with $m\leq n$
and a matrix $M:\Rn\rightarrow\Rm$ with $Z$ non-zeros, let $\mathcal{T}$
be the cost to compute $\nabla\delta_{K_{1}}^{*}$ and $\nabla\delta_{K_{2}}^{*}$.
Then, we can solve the problem
\[
\min_{x\in K_{1},Mx\in K_{2}}c^{\top}x
\]
in time $O^{*}(m\mathcal{T}+mZ+m^{3}).$
\end{thm}

\begin{rem*}
We hid all sorts of terms in the log term hidden in $O^{*}$ such
as the diameter of the set. Also this is the number of arithmetic
operations, not the bit complexity.
\end{rem*}
Going back to the school/student problem, this algorithm gives a running
time of
\[
O^{*}(|V_{2}||E|+|V_{2}|^{3})
\]
which is linear in the number of students!

In general, this statement says that if we can split a convex problem
into two parts, with both being easy to solve and one part having
fewer variables, then we can solve the entire problem in time depending
on the smaller dimension.
\begin{xca}
How fast can we solve $\min_{x\in\cap_{i=1}^{k}K_{i}}c^{\top}x$ given
the oracles $\nabla\delta_{K_{i}}^{*}$?
\end{xca}


